\section{Related Work}
\label{sec:related}

\para{Single-model collapse.}
\citet{shumailov2024nature} formalised model collapse, showing that a
single OPT-125m fine-tuned on its own \wikitext outputs has perplexity
that rises monotonically across generations and a generation
distribution whose tails progressively disappear.
\citet{alemohammad2023mad} extended the result to image generative
models under the name \emph{Model Autophagy Disorder} (MAD), and
\citet{briesch2023suffer} disentangled quality decay (preserved on
verifiable tasks) from diversity decay (not preserved).
\citet{guo2023linguistic} introduced lexical, syntactic, and semantic
diversity metrics for the recursive-training setting; we reuse the
distinct-bigram metric from that work. Theoretical accounts include
\citet{dohmatob2024tale}, who frame collapse as a change of scaling laws,
and \citet{seddik2024synthetic}, who derive a maximal admissible
synthetic ratio. Our N=1 condition is a direct reproduction of the
\citet{shumailov2024nature} ``replace'' regime.

\para{Mitigation by data dynamics.}
\citet{gerstgrasser2024accumulating} show that \emph{accumulating}
synthetic and real data (rather than replacing) yields a finite,
bounded test error independent of iteration count, both empirically
on TinyStories and theoretically. \citet{dey2024universality} sharpen
this to a universal $\pi^2/6$ asymptote. These results establish
data-side mitigations; we deliberately ablate them by running our
fine-tuning harness in the harshest replace regime so that diversity
alone is the variable being tested.

\para{Ecosystem-level collapse.}
The work most directly comparable to ours is
\citet{hodel2026epistemic}: $M$ copies of one base model, each
fine-tuned on a non-overlapping segment of \wikitext, then jointly
generating, pooling, shuffling, and re-splitting the synthetic data
each iteration. They sweep $M \in \{1, 2, 4, 16\}$ and find the
\emph{optimal} $M$ grows monotonically with iteration count. Our E1
reproduces their compute-controlled setup but (a) extends the sweep
to $N{=}8$ at finer spacing, (b) reports an explicit \mvp estimate at
fixed horizon $T{=}6$, and (c) provides a direct comparison to the
much cheaper retrieval-based ecosystem.

\citet{wang2025web} build a parameter-frozen retrieval ecosystem of
$N{=}3$ pretrained-different LLMs sharing a growing post pool, and
show that the Frobenius norm of the inter-agent embedding-distance
matrix converges to a small constant. Our E2 generalises their setup
along the missing $N$ axis. \citet{vu2025recursive} provide a
theoretical framework for $N{=}2$ heterogeneous models, validate it
on Llama and Phi family LLMs, and document the substantial overlap of
modern LLM pretraining corpora --- empirical motivation for the
ecosystem framing.

\citet{kovac2025recursive} treat the data side as the diversity knob,
rotating across four small base models per generation and studying
how training-data properties (lexical, semantic, quality) modulate
distribution shift. Their result that lexical diversity \emph{amplifies}
shift while semantic diversity \emph{mitigates} it is consistent with
our finding that prompt-induced ``surface'' diversity buys little
collapse-resistance, while pretraining-prior diversity buys a lot.

\para{Population ecology and minimum viable population.}
\citet{shaffer1981minimum} introduced the concept of \mvp as the
smallest isolated population with a high probability of long-run
persistence in the wild. \citet{traill2007mvp} report a generic
vertebrate \mvp around several thousand individuals; our analogue is
specifically about the \emph{number of distinct LLMs} required to keep
their shared synthetic-text substrate from converging. The Hill--Shannon
Diversity / Vendi score \citep{friedman2023vendi} we use as a primary
metric is itself imported from ecology.

\para{Positioning.}
We are not the first to observe that ecosystem-level diversity helps;
\citet{hodel2026epistemic} and \citet{wang2025web} both make this
qualitative point. Our contribution is to (i) give the first explicit
\mvp estimates that pin down a number, (ii) compare four axes of
``individuality'' on identical metrics, and (iii) cross-check between
the two methodologically very different ecosystem operationalisations
that until now had not been benchmarked side by side.
